% Anonymous LASS 2026 submission draft.
% Workshop requirements: ACM CIKM 2026 two-column format, 4--8 pages
% excluding references, and double-blind review.
\documentclass[sigconf,anonymous]{acmart}

\usepackage{amsmath}
\usepackage{tabularx}

% Submission copy: suppress metadata assigned only after acceptance.
\setcopyright{none}
\settopmatter{printacmref=true}
\renewcommand\footnotetextcopyrightpermission[1]{}

\newcommand{\method}{\textsc{ScopeProbe}}
\newcommand{\egm}{\textsc{Evidence-Gated Memory}}
\newcommand{\ourscope}{\mathcal{Z}}

\begin{document}

\title{The Right Event, the Wrong Lesson: Diagnosing Reflective
Misconsolidation in Social LLM Agents}

% Keep the source anonymous during review.
\author{Anonymous Author(s)}
\affiliation{%
  \institution{Anonymous Institution}
  \country{Anonymous}
}

\begin{abstract}
Long-running social simulations require language-model agents to convert
interaction histories into persistent beliefs about themselves, particular
partners, and the shared environment. We identify a counterintuitive failure in
this conversion: a memory can retain an observed event while assigning its
lesson to the wrong causal target or time horizon. We call this
\emph{reflective misconsolidation}. We introduce \method, an
intervention-based diagnostic harness that constructs matched social
trajectories with similar outcomes but different latent causes, then jointly
audits persistent memory, an elicited attribution readout, and later action.
Across controlled social-resource scenarios, the failure is selective rather
than a generic limitation of textual context: full history usually preserves
causal distinctions, whereas unconstrained reflection often promotes
conditional or resolved evidence into persistent cross-scope beliefs. Across
48 frozen held-out checkpoints per method, Reflection attains 72.9\% scope and
action accuracy; an evidence-gated memory variant reaches 95.8\% on both under
the same decision prompt. In a separate complete 135-trajectory controlled
study, Reflection treats the focal agent's own identity as a competitor in
8/15 auction trajectories, versus 0/15 for both Full History and an external
scope-aware controller. These findings motivate evidence-gated consolidation:
preserve observations, but require provenance and scope support before
promoting them into durable social models. Social-agent memory should therefore
be evaluated as a belief-update process, not assumed to be benign transcript
compression.
\end{abstract}

\begin{CCSXML}
<ccs2012>
 <concept>
  <concept_id>10010147.10010257</concept_id>
  <concept_desc>Computing methodologies~Machine learning</concept_desc>
  <concept_significance>500</concept_significance>
 </concept>
 <concept>
  <concept_id>10010405.10010455</concept_id>
  <concept_desc>Applied computing~Sociology</concept_desc>
  <concept_significance>300</concept_significance>
 </concept>
</ccs2012>
\end{CCSXML}

\ccsdesc[500]{Computing methodologies~Machine learning}
\ccsdesc[300]{Applied computing~Sociology}

\keywords{LLM agents, social simulation, agent memory, reflection,
causal attribution, evaluation}

\acmConference[LASS '26]{The 2nd Workshop on LLM Agents for Social
Simulation}{November 8, 2026}{Rome, Italy}
\acmYear{2026}
\acmDOI{}
\acmISBN{}

\begin{teaserfigure}
  \centering
  \includegraphics[width=\textwidth]{figures/scopeprobe_overview.pdf}
  \caption{The same adverse outcome can license different persistent beliefs.
  A world-level capacity change and a named participant's temporary emergency
  should not be consolidated into the same global rule. \method{} audits the
  resulting event $\rightarrow$ memory $\rightarrow$ attribution
  $\rightarrow$ later-action chain, while evidence gating routes claims to
  their supported causal target and temporal status. A concrete identity
  failure appears in our controlled auction: although the role says ``You are
  Erin,'' Reflection reasons that it must ``bid strictly above Erin'' in 8/15
  trajectories (Full History: 0/15; scope-aware controller: 0/15).}
  \Description{Four panels show matched social-resource events with different
  causes, their collapse into a generic belief under reflection, an audit chain
  from observation through memory and attribution to later action, and an
  evidence gate routing claims to self, named-other, world, or temporal-status
  channels. A lower evidence strip shows a focal agent named Erin mistakenly
  treating Erin as a competing bidder.}
  \label{fig:overview}
\end{teaserfigure}

\maketitle

\section{Introduction}

In a long-running social simulation, memory is not merely a convenience for
recovering old dialogue; it is part of the simulator's evolving social state.
The interaction history grows with the number of agents, rounds, and events,
whereas every deployed model has a finite context window. Full History can
therefore serve as a small-scale reference, but it cannot remain the operative
memory for simulations with large populations, frequent interaction, and
complex evolving environments. Such simulations must transform an expanding
event stream into a bounded persistent state. Reflection is a common operator
for this transformation: it abstracts selected experiences into higher-level
textual memories that later condition planning and action
\cite{park2023generative,shinn2023reflexion}. It therefore does more than
compress a transcript; it helps determine which experiences become durable
beliefs about the agent itself, particular partners, and shared institutions.
Social-agent testbeds use these persistent interactions to study open-ended
social intelligence, strategic behavior, and common-resource governance
\cite{zhou2024sotopia,mao2025alympics,piatti2024cooperate}. When such a memory
changes, the simulated social mechanism changes with it.

Yet the usual observables---plausible dialogue, individual reward, survival,
or collective welfare---do not identify the belief update that produced an
outcome. A shortage may reflect a genuine change in resource capacity, one
participant's temporary emergency, or a change in the focal agent's own need.
Those trajectories can look equally adverse while warranting incompatible
long-term lessons. This gap matters if agent simulations are to function as
scientific instruments: role-play plausibility and aggregate behavior do not
by themselves make the internal mechanism faithful or auditable
\cite{li2026agentsalone}.

We expose a failure hidden by outcome-only evaluation: \emph{the event remains
legible, but its epistemic address changes}. As illustrated in
Figure~\ref{fig:overview}, a reflection can retain that a named participant
overused a resource yet rewrite the event as evidence that the shared system is
globally unstable. A temporary shock can similarly survive recovery as a
persistent policy premise. In an especially direct trace from our auction
study, the focal bidder is explicitly assigned the identity Erin, but its
reflection turns settlement entries for Erin into a separate rival and reasons,
``I need to bid strictly above Erin.''

We call the erroneous promotion of evidence into an unsupported durable claim \emph{reflective
misconsolidation}; its cross-target and cross-time manifestations are
\emph{scope interference}. Unlike a retrieval failure, the relevant event is
still available. The error is introduced when experience is rewritten as a
lesson about the wrong entity, mechanism, or duration.

To isolate this process, we introduce \method, a diagnostic harness built from
matched social trajectories. Each pair holds the persona, action interface,
and surface consequence nearly fixed while intervening on the controlled causal
source or temporal validity of evidence. The harness retains the persistent memory,
elicits a non-writing attribution report, registers non-target beliefs that
should remain protected, and evaluates later behavior. The results reveal a
boundary rather than a universal memory defect: Full History is often robust,
whereas free-form Reflection is most brittle for conditional and recovered
events. Table~\ref{tab:heldout} shows the frozen held-out diagnostic. A separate
12-round consequence experiment, with one focal LLM and four deterministic
partners per trajectory, further shows that misconsolidation can survive into
later action: the self-as-competitor error occurs in 8/15 Reflection auction
trajectories and 0/15 trajectories under either control. Guided by the
diagnosis, we instantiate a lightweight mitigation, \egm, that separates
observations from persistent models, gates consolidation by evidence and
causal target, and keeps immediate reversible precautions distinct from
long-term policy. Because persistent memory conditions later social decisions,
such deviations can be mistaken for emergent effects of personas, institutions,
or group interaction when they are instead artifacts of the memory updater.

Our contributions are threefold:
\begin{itemize}
  \item \textbf{A diagnostic problem formulation.} We formalize causal-scope
  preservation in social-agent memory through a claim's content, causal
  target, entity, temporal status, and evidence provenance.
  \item \textbf{A matched diagnostic.}\\[-1pt]
  \method{} traces event $\rightarrow$ persistent memory $\rightarrow$ elicited
  attribution $\rightarrow$ later action under matched causal-source
  interventions and protected-belief controls.
  \item \textbf{Mechanistic findings and a diagnosis-derived mitigation.} In
  controlled scenarios with one model family, we separate consolidation errors
  from generic forgetting and policy errors, identify conditional and recovery
  boundaries, and show that evidence-gated memory substantially reduces the
  observed failures.
\end{itemize}

The resulting takeaway is methodological: social simulation should ask not
only whether agents reach a plausible outcome, but whether they learn the
right social model for the right reason.

\begin{table}[t]
  \caption{Frozen held-out diagnostic: 48 clustered checkpoints per row
  (8 scenarios $\times$ 3 repeats $\times$ 2 checkpoints). Leakage is mean
  probe-reported update strength on registered non-target labels, not a claim
  error rate; action correctness uses registered intervals. Reflection and
  \egm{} share a decision prompt.}
  \label{tab:heldout}
  \centering
  \small
  \setlength{\tabcolsep}{3.4pt}
  \begin{tabular}{@{}lccc@{}}
    \toprule
    Memory updater & Scope $\uparrow$ & Leakage $\downarrow$ & Action $\uparrow$ \\
    \midrule
    Reflection            & 72.9 & 0.0979          & 72.9 \\
    Guarded Reflection    & 91.7 & \textbf{0.0028} & 75.0 \\
    Structured Reflection & 91.7 & 0.0398          & 87.5 \\
    \egm                  & \textbf{95.8} & \textbf{0.0028} & \textbf{95.8} \\
    \bottomrule
  \end{tabular}
\end{table}

\section{Related Work}

\subsection{LLM Agents for Social Simulation}

Generative Agents established a practical architecture in which agents retrieve
selected observations, synthesize higher-level reflections, and use both to
plan later behavior \cite{park2023generative}. Rather than placing an ever-growing
transcript into every decision context, this design makes reflection a
consolidation operator between experience and subsequent action.
Subsequent environments broadened the evaluation target: SOTOPIA studies social
goal pursuit in interactive scenarios \cite{zhou2024sotopia}; ALYMPICS uses
language agents for controlled strategic games \cite{mao2025alympics}; and
GovSim studies cooperation and sustainability in common-resource societies
\cite{piatti2024cooperate}. Controlled agent populations have also been used to
probe social mechanisms such as the interaction between history, persona, and
spiral-of-silence dynamics \cite{zhong2025spiral}. These works demonstrate the
scientific potential of agent simulation, but their primary readouts are
interaction quality, strategy, or population outcome. \method{} is
complementary: it treats the persistent belief update that generates those
outcomes as the object of intervention. Our resource scenarios are synthetic
diagnostic adaptations inspired by GovSim and ALYMPICS, not executions of the
original closed-loop benchmarks.

\subsection{Reflection and Persistent Agent Memory}

Reflection abstracts selected experience into higher-level text in Generative
Agents \cite{park2023generative}, while Reflexion stores verbal feedback to
improve later trials \cite{shinn2023reflexion}. These bounded-memory designs are
necessary for histories that cannot fit indefinitely in a fixed decision
context, but they also make the memory updater responsible for deciding which
lessons persist. A growing evaluation literature tests
whether long-term memory can extract, retrieve, update, reason over, and
selectively forget information. LongMemEval emphasizes multi-session and
temporal reasoning \cite{wu2025longmemeval}; MemoryAgentBench evaluates four
interactive competencies including selective forgetting
\cite{hu2026memoryagentbench}; and HaluMem localizes hallucinations to
extraction, update, or question-answering operations \cite{chen2026halumem}.
Empirical work further shows that stored experiences can propagate errors or
induce misaligned replay \cite{xiong2026memory}, while TrustMem explicitly
optimizes consolidation for coverage, preservation, and faithfulness
\cite{yang2026trustmem}. Concurrent evidence shows that compression can erase
an epistemic qualifier even when the underlying claim is retained
\cite{kwon2026explicit}. These studies establish that memory can be unfaithful,
stale, or lossy. Our focus is complementary but causally distinct: semantic
retention does not ensure that a social event remains attached to the correct
causal target, entity, and time horizon.

\subsection{Attribution, Temporal Validity, and Memory-to-Action}

The closest work exposes individual facets of this problem. PASB identifies
status promotion, attribution removal, and scope broadening when user claims
are committed to a personal agent's durable state \cite{mao2026pasb}. STALE
tests implicit invalidation of old memories and downstream policy adaptation
\cite{chao2026stale}. Actor--observer studies reveal perspective-sensitive
failure attribution in reflection and external auditing
\cite{li2026actorobserver}. Finally, Mem2ActBench and MemoryArena connect
long-term memory to tool use and interdependent later actions
\cite{shen2026mem2act,he2026memoryarena}. We do not claim priority over scope
broadening, stale memory, attribution bias, or memory-conditioned action.
Instead, \method{} combines them in a controlled social-agent setting: matched
trajectories vary the latent cause while holding the visible consequence
similar, factor attribution into self, named-other, world, and temporal status,
measure leakage into registered protected beliefs, and trace the result into
later social action. Existing work asks whether memory is faithful, current,
or useful; we ask whether social evidence is consolidated at the correct
\emph{epistemic address}.

\section{Preliminaries and Problem Formulation}
\label{sec:prelim}

\subsection{Social Agent with Persistent Memory}

Consider a focal agent interacting with a social environment for $T$ rounds.
At round $t$, it receives observation $o_t$, acts using the persistent memory
available before the current write, and observes an outcome $y_t$ produced
jointly by its action, the other agents' actions $a_t^{-}$, and environment
state $s_t$. It then updates that memory:
\begin{equation}
  \begin{aligned}
    a_t &\sim \pi(\cdot \mid o_t,m_{t-1}),\\
    (s_{t+1},y_t)&\sim\mathcal{T}(s_t,a_t,a_t^{-}),\\
    m_t &= U(m_{t-1},o_t,a_t,y_t).
  \end{aligned}
  \label{eq:agent-loop}
\end{equation}
Here $U$ may append full history, summarize, reflect, or invoke an external
memory controller. A measurement query $Q$ produces an attribution readout
$\hat b_t=Q(m_t)$ without writing the query or response back to memory. Thus the
causal sequence under study is
\begin{equation}
  (o_t,a_t,y_t) \longrightarrow m_t \longrightarrow \hat b_t
  \longrightarrow a_{t+1:T},
  \label{eq:chain}
\end{equation}
not a claim that the newly written memory caused the already chosen $a_t$.
The readout is an observable self-report, not direct access to a latent belief.

\subsection{Epistemic Address}

Let a memory claim be
\begin{equation}
  c=(\phi,z,e,\tau,w),
  \label{eq:claim}
\end{equation}
where $\phi$ is semantic content, $z\in\ourscope$ is its causal target,
$e$ is the referenced entity, $\tau$ is temporal status, and $w$ records the
supporting evidence or provenance. We use four causal targets---\textsc{self},
\textsc{named-other}, \textsc{world}, and \textsc{persona}---and factor time
separately as \textsc{temporary}, \textsc{persistent}, or \textsc{resolved}.
The triple
\begin{equation}
  \alpha(c)=(z,e,\tau)
  \label{eq:address}
\end{equation}
is the claim's \emph{epistemic address}: whom or what the claim concerns and
for how long it is licensed. Remembering content $\phi$ is therefore
insufficient. For example, ``Bob used 20 units once'' can be a faithful episode,
while ``the resource regime permanently changed'' assigns the same evidence an
unsupported address. Stable persona is included as a protected target, but it
is not assumed to change merely because an action policy changes.

Our operational probe uses a compact label space that mixes these factored
dimensions: \texttt{self}, \texttt{other}, \texttt{world}, and
\texttt{persona} denote causal targets; \texttt{episodic} denotes a temporary
event; and \texttt{none} denotes a resolved event or no licensed persistent
update. The factorization in Equation~\ref{eq:address} makes this implementation
choice explicit rather than treating all six labels as equivalent ontological
scopes.

\subsection{Reflective Misconsolidation}

Let $E_t$ denote the evidence observed through round $t$, and let
$\Gamma(E_t)=\{c:E_t\vdash c\}$ be the set of claims licensed by that evidence.
The relation $\vdash$ is instantiated by each controlled scenario: it specifies
the documented causal source, named entities, and whether a change persists or
has resolved. Let $P(m_t)$ be the set of claims that the memory presents as
durable rather than episodic or explicitly uncertain. Misconsolidation occurs
when
\begin{equation}
  \mathcal{M}_t=P(m_t)\setminus\Gamma(E_t)\neq\emptyset.
  \label{eq:misconsolidation}
\end{equation}
When $U$ is a reflective updater, we call this \emph{reflective
misconsolidation}. \emph{Scope interference} is the observable subset in which
a claim is grounded in an observed event but its address $\alpha(c)$ is not
licensed. It includes: (i) cross-target misattribution, such as named-other
evidence promoted to a world rule; (ii) entity overgeneralization, such as one
participant promoted to the group; (iii) temporal overextension, such as a
resolved anomaly retained as persistent; and (iv) mechanism invention, where a
surface outcome is assigned an unsupported hidden cause. Retaining a single
event as an episode or uncertain hypothesis is not an error. Nor does weak
evidence for persistence preclude a minimal, reversible response to a verified
current hazard.

\subsection{Matched-Intervention Diagnostic Objective}

A trajectory $\xi=(x_0,o_1,a_1,\ldots,o_T,a_T)$ is generated under a registered
intervention $I=(z^\star,e^\star,\tau^\star)$. \method{} pairs trajectories
whose personas, interfaces, and visible adverse outcomes are similar, while
changing the intervention's latent causal target or temporal status. This
prevents the outcome alone from revealing which belief should change.

At checkpoint $t$, the non-writing probe returns a primary label $\hat y_t$
and update strengths $r_t(k)\in[0,1]$ over operational labels $k$. Let
$g(z^\star,\tau^\star)$ map the factored target to its registered probe label,
and let $\mathcal{P}_t$ be labels that should remain unchanged. The principal
diagnostics are
\begin{align}
  \mathrm{ScopeAcc}_t
    &=\mathbf{1}\!\left[\hat y_t=g(z^\star,\tau^\star)\right],
    \label{eq:scope-acc}\\
  \mathrm{Leakage}_t
    &=\frac{1}{|\mathcal{P}_t|}\sum_{k\in\mathcal{P}_t} r_t(k),
    \label{eq:leakage}\\
  \mathrm{Margin}_t
    &=r_t\!\left(g(z^\star,\tau^\star)\right)
      -\max_{k\in\mathcal{P}_t}r_t(k).
    \label{eq:margin}
\end{align}
Leakage is therefore a probe-reported protected-update strength, not a claim
error rate. We separately inspect persistent text and score whether actions
fall inside scenario-registered acceptable sets. A memory-to-action consequence
is attributed only to actions chosen after the relevant write, as in
Equation~\ref{eq:chain}. The diagnostic objective is not simply to maximize
recall: it is to strengthen the licensed target while preserving protected
beliefs and producing a proportionate downstream response.

% Method, experiments, limitations, and conclusion will be integrated by the
% paper team after this complete front-half draft.

\bibliographystyle{ACM-Reference-Format}
\bibliography{references}

\end{document}
